% ============================================================================================
% Appendices. App.~A holds definitions omitted from Chapters 2 and 3; App.~B holds supplementary
% tables. Numbers are transcribed from src/features/* and src/analysis/*.
% ============================================================================================

\chapter{Mathematical Definitions}
\label{app:formulas}
Chapters~\ref{ch:background} and~\ref{ch:methodology} define the tracking metric, the four distances and
the regression in the body. This appendix supplies only what was left out of them: the scene covariates, the design matrix, and the two
statistics used to report fit and out-of-sample gain.

\section{Scene covariates}
\label{app:covariates}
Three interpretable covariates accompany the environment distance (Section~\ref{sec:meth-features}). Writing
$f$ for a frame, $\ell$ for a location, and the bar for a mean over that location's frames,
\begin{align}
  \text{achromatic}(f) &= \Big[\; \overline{\max_c f_c - \min_c f_c} \;<\; \tau_{\mathrm{ir}}\;\Big],\quad \tau_{\mathrm{ir}}=12,\\
  \mathrm{is\_night\_ir}(\ell) &= \big[\; \overline{\text{achromatic}(f)} > 0.5 \;\big],\qquad
  \mathrm{clutter}(\ell) = \overline{\operatorname{Var}\!\big(\nabla^2 f\big)},
\end{align}
where $\max_c-\min_c$ is the per-pixel spread across colour channels and $\nabla^2$ is the four-neighbour
Laplacian on the greyscale background.

\section{The design matrix}
\label{app:design}
Equation~\ref{eq:glm} fits $\operatorname{logit}\mathbb{E}[y]=x^{\!\top}\beta$ with variance weights
$w=n_{\mathrm{frames}}$. The design $x$ standardises each continuous predictor as $z=(x-\mu)/\sigma$, imputing
missing values to the training mean $\mu$ and dropping zero-variance columns. It adds the standardised
$\ln(1+n_{\mathrm{frames}})$ as the support covariate, passes the binary $\mathrm{is\_night\_ir}$ through
unchanged, and adds an intercept. Every standardisation statistic is learned on training rows only.

\section{Fit and out-of-sample gain}
\label{app:stats}
Fit quality is the McFadden pseudo-$R^2$. The per-cell out-of-sample gain over a mean predictor is $\delta_i$:
\begin{equation}
  R^2_{\mathrm{McF}}=1-\frac{D}{D_0},
  \qquad
  \delta_i=\underbrace{|y_i-\bar y|}_{\text{baseline error}}-\underbrace{|y_i-\hat y_i|}_{\text{model error}},
\end{equation}
with $D$ and $D_0$ the model and null (intercept-only) deviances. A positive $\delta_i$ means the distance
model has lower absolute error than the mean predictor on cell $i$.

\paragraph{Planted-effect power.} For a planted standardised effect of size $\beta$ on a feature $z$, each
replicate transforms the observed score $y$, clipped to $[\varepsilon, 1-\varepsilon]$ with
$\varepsilon=10^{-3}$, as
\begin{equation}
  y' \;=\; \sigma\big(\operatorname{logit}(y) + \beta z\big),
  \qquad
  \operatorname{logit}(y)=\ln\frac{y}{1-y},
  \quad
  \sigma(u)=\frac{1}{1+e^{-u}},
  \label{eq:plant}
\end{equation}
and the unchanged out-of-sample procedure of Section~\ref{sec:meth-cv} is asked to detect the plant. Power is
the detected fraction over replicates (Table~\ref{tab:mde}).

For the species-level false-positive analysis (Section~\ref{sec:meth-fp}), correlations are weighted by the
per-species probe counts $w$:
\begin{equation}
  r_w=\frac{\sum_i w_i (x_i-\bar x_w)(y_i-\bar y_w)}
           {\sqrt{\sum_i w_i (x_i-\bar x_w)^2}\,\sqrt{\sum_i w_i (y_i-\bar y_w)^2}},
  \qquad \bar x_w=\frac{\sum_i w_i x_i}{\sum_i w_i}.
\end{equation}

\chapter{Supplementary Tables}
\label{app:supplementary}
These tables support claims made in Chapter~\ref{ch:results}. Each is referenced from the section that
discusses it.

\IfFileExists{tables/variance.tex}{% generated by src/analysis/report.py --- do not edit by hand
\begin{table}[htbp]\centering
  \caption{Variance partition for detection (\texttt{pDetA}): each distance's unique share of $R^2$ (LMG/Shapley) and its VIF. The four distances together explain only $\approx3\%$ of the variance, no factor dominates, and every VIF~$\approx1$ --- the null is not masked by collinearity.}
  \label{tab:variance}
  \begin{tabular}{lccc}
    \toprule
    Factor & unique $R^2$ & share & VIF \\
    \midrule
    Taxonomic & 0.014 & 48\% & 1.06 \\
    Visual & 0.012 & 39\% & 1.06 \\
    Environment & 0.004 & 13\% & 1.00 \\
    Temporal & 0.000 & 0\% & -- \\
    \bottomrule
  \end{tabular}
\end{table}
}{}
\IfFileExists{tables/difficulty.tex}{% generated by src/analysis/report.py --- do not edit by hand
\begin{table}[htbp]\centering
  \caption{Nested size decomposition of the intrinsic-difficulty test (detection, 346 cells). Each entry is the out-of-sample MAE gain over a mean-predictor baseline, $\Delta$~[95\% CI]~($p$), under leave-species-out and leave-location-out. The entire gain sits in the size covariate ($\log$-area); low-light \emph{alone} does not validate --- the difficulty signals only pass the bar by riding on size.}
  \label{tab:difficulty}
  \begin{tabular}{lcc}
    \toprule
    Model & Leave-species-out $\Delta$~[CI]~($p$) & Leave-location-out $\Delta$~[CI]~($p$) \\
    \midrule
    $\log$(support) only & +0.001~[-0.007,\,+0.008]~(0.432) & +0.002~[-0.005,\,+0.008]~(0.291) \\
    $\log$(area) only \emph{[size]} & +0.015~[+0.002,\,+0.028]~(0.014) & +0.017~[+0.009,\,+0.025]~(0.001) \\
    Low-light only (no size) & +0.003~[-0.006,\,+0.011]~(0.233) & +0.004~[-0.003,\,+0.011]~(0.123) \\
    Low-light $+$ size & +0.017~[+0.004,\,+0.032]~(0.011) & +0.020~[+0.011,\,+0.030]~(0.000) \\
    \bottomrule
  \end{tabular}
\end{table}
}{}
\IfFileExists{tables/gate1_measurement_species.tex}{% generated by src/analysis/report.py --- do not edit by hand
\begin{table}[t]\centering
  \caption{Zero-shot SAM\,3 on the pooled train+test set --- support-weighted mean over 1025 scored cells (of 2126).}
  \label{tab:gate1_species}
  \begin{tabular}{lr}
    \toprule
    Metric (mask) & Value \\
    \midrule
    pDetA & 0.536 \\
    pAssA & 0.597 \\
    pHOTA & 0.559 \\
    \bottomrule
  \end{tabular}
\end{table}
}{}
\IfFileExists{tables/gate1_measurement_burst.tex}{% BURST Gate-1 --- hand-maintained (mask-HOTA); see docs/burst_replication.md
\begin{table}[htbp]\centering
  \caption{Zero-shot SAM\,3 on the BURST val animal subset --- support-weighted mean over 132 probe cells
  (41 species), mask-HOTA. The score is non-degenerate and the association target is genuine
  ($\mathrm{corr}(\texttt{pDetA},\texttt{pAssA})=0.76$, $71/132$ multi-object cells).}
  \label{tab:gate1_burst}
  \begin{tabular}{lr}
    \toprule
    Metric (mask-HOTA) & Value \\
    \midrule
    pDetA & 0.607 \\
    pAssA & 0.688 \\
    pHOTA & 0.635 \\
    \bottomrule
  \end{tabular}
\end{table}
}{}
\IfFileExists{tables/cv_validation_species.tex}{% generated by src/analysis/report.py --- do not edit by hand
\begin{table}[htbp]\centering
  \caption[Out-of-sample validation, species hold-out]{Out-of-sample error vs a mean predictor on the \textbf{species hold-out} (Split~A; the $1025$ pooled
  train+test cells over $99$ species). Each row is one cross-validation grouping scheme --- leave-species-out or
  leave-location-out --- applied within these same cells; $\Delta$ is the MAE gain over baseline, with a
  paired-bootstrap 95\% CI and one-sided $p$. No scheme's gain is significant. (Distinct experiment from
  Table~\ref{tab:cv}, hence the different $n$ and values.)}
  \label{tab:cv_species}
  \begin{tabular}{llrrrcr}
    \toprule
    CV scheme & Target & $n$ & MAE & Baseline & $\Delta$~[95\% CI] & $p$ \\
    \midrule
    location & pAssA & 1025 & 0.270 & 0.270 & -0.000~[-0.006,\,+0.005] & 0.52 \\
    location & pDetA & 1025 & 0.274 & 0.274 & -0.001~[-0.007,\,+0.006] & 0.58 \\
    species & pAssA & 1025 & 0.273 & 0.270 & -0.003~[-0.014,\,+0.008] & 0.74 \\
    species & pDetA & 1025 & 0.277 & 0.274 & -0.004~[-0.015,\,+0.009] & 0.74 \\
    \bottomrule
  \end{tabular}
\end{table}
}{}
\IfFileExists{tables/ablations.tex}{% generated by src/analysis/report.py --- do not edit by hand
\begin{table}[htbp]\centering
  \caption{Factor ablation for detection (\texttt{pDetA}): McFadden pseudo-$R^2$ and the out-of-sample MAE gain over a mean baseline under leave-species-out (LSO) and leave-location-out (LLO). Dropping any single distance barely changes the fit or the (null) out-of-sample gain --- no factor carries signal; only adding the animal-size covariate moves the score.}
  \label{tab:ablation}
  \begin{tabular}{lccc}
    \toprule
    Design & pseudo-$R^2$ & LSO $\Delta$~($p$) & LLO $\Delta$~($p$) \\
    \midrule
    full (4 distances) & 0.055 & $+0.004$~($p{=}0.24$) & $+0.004$~($p{=}0.23$) \\
    drop taxonomic & 0.039 & $+0.000$~($p{=}0.47$) & $+0.000$~($p{=}0.49$) \\
    drop temporal & 0.055 & $+0.004$~($p{=}0.24$) & $+0.004$~($p{=}0.23$) \\
    drop visual & 0.048 & $+0.004$~($p{=}0.22$) & $+0.004$~($p{=}0.28$) \\
    drop environment & 0.055 & $+0.006$~($p{=}0.15$) & $+0.005$~($p{=}0.17$) \\
    $+$ animal size & 0.119 & $+0.018$~($p{=}0.02$) & $+0.020$~($p{=}0.00$) \\
    $-$ support covariate & 0.035 & $-0.002$~($p{=}0.72$) & $-0.003$~($p{=}0.73$) \\
    \bottomrule
  \end{tabular}
\end{table}
}{}
\IfFileExists{tables/robustness.tex}{% generated by src/analysis/report.py --- do not edit by hand
\begin{table}[htbp]\centering
  \caption{Design-robustness sweep (detection). Each row re-runs the primary model with one design choice swapped; entries are the out-of-sample MAE gain, $\Delta$~[95\% CI]~($p$), under leave-species-out (LSO) and leave-location-out (LLO) (Bonferroni $m=5$). The null holds throughout --- no encoder, crop, prompt, or distributional distance beats the mean.}
  \label{tab:robustness}
  \footnotesize
  \setlength{\tabcolsep}{2.5pt}
  \begin{tabular}{>{\raggedright\arraybackslash}p{2.9cm}llc}
    \toprule
    Variant & LSO $\Delta$~[CI]~($p$) & LLO $\Delta$~[CI]~($p$) & Beats mean? \\
    \midrule
    baseline (DINOv2, mask-crop, species prompt) & +0.004~[-0.006,\,+0.015]~(0.237) & +0.004~[-0.007,\,+0.015]~(0.235) & n.s. \\
    CLIP encoder & +0.004~[-0.006,\,+0.014]~(0.247) & +0.004~[-0.007,\,+0.014]~(0.270) & n.s. \\
    whole-frame (bbox + background) & +0.002~[-0.008,\,+0.013]~(0.349) & +0.002~[-0.008,\,+0.013]~(0.357) & n.s. \\
    generic ``animal'' prompt & -0.346~[-0.381,\,-0.309]~(1.000) & -0.343~[-0.363,\,-0.326]~(1.000) & n.s. \\
    Frechet distributional distance & +0.007~[-0.005,\,+0.019]~(0.120) & +0.007~[-0.004,\,+0.019]~(0.102) & n.s. \\
    MMD distributional distance & +0.002~[-0.008,\,+0.012]~(0.374) & +0.002~[-0.009,\,+0.012]~(0.349) & n.s. \\
    \bottomrule
  \end{tabular}
\end{table}
}{}
\IfFileExists{tables/familiarity.tex}{% SAM 3 familiarity proxy (T2.5) --- hand-maintained; reproduce with src/analysis/familiarity_experiment.py
\begin{table}[htbp]\centering
  \caption{SAM\,3 familiarity proxy (T2.5, detection): does the separability of a species in SAM\,3's \emph{own}
  vision-encoder feature space predict \texttt{pDetA}? Each metric is size-controlled and validated at the same
  leave-species-out bar as the distances; $\Delta$MAE~($p$) is the out-of-sample gain over a mean predictor.
  None clears the Bonferroni threshold ($\alpha=0.05/3=0.017$) on the leave-species-out (novelty) test on
  either dataset. The last columns (SA-FARI) are the cell-level correlations with the DINOv2 visual distance and
  animal size: all three metrics largely \emph{re-encode} the visual distance ($r_\text{vis}=0.63$--$0.85$).
  On the same $346$ SA-FARI cells the before-running animal-size control validates out-of-sample, so this is a
  genuine null, not an underpowered one.}
  \label{tab:familiarity}
  \begin{tabular}{lccccc}
    \toprule
    Metric & SA-FARI $\Delta$MAE~($p$) & BURST $\Delta$MAE~($p$) & $r_\text{vis}$ & $r_\text{size}$ & clears bar? \\
    \midrule
    silhouette        & +0.012~(0.034) & +0.012~(0.020) & +0.67 & $-$0.16 & no \\
    nearest-prototype & +0.009~(0.068) & +0.006~(0.136) & +0.63 & +0.03 & no \\
    mahalanobis       & +0.008~(0.085) & +0.006~(0.148) & +0.85 & +0.32 & no \\
    \bottomrule
  \end{tabular}
\end{table}
}{}
\IfFileExists{tables/model_swap_cv.tex}{\input{tables/model_swap_cv}}{}
\IfFileExists{tables/replication_burst.tex}{% BURST replication --- hand-maintained; reproduce with scripts/burst_replication_analysis.py
\begin{table}[htbp]\centering
  \caption{BURST independent replication --- out-of-sample error under leakage-free leave-species-out
  cross-validation (41 species, 132 mask-native cells), against a mean predictor. $\Delta$ is the MAE gain over
  baseline for \texttt{pDetA}; the bracketed 95\% CI and one-sided $p$ come from a paired bootstrap resampling
  whole held-out species. Every before-running distance is null (CI spans zero); not even the animal-size
  control reaches significance here, so BURST is a convergent null that cannot independently certify power for a
  small species-level effect. BURST has no shared locations, so leave-species-out is the only valid scheme.}
  \label{tab:replication_burst}
  \begin{tabular}{lrcr}
    \toprule
    Predictor & $\Delta$MAE & 95\% CI & $p$ \\
    \midrule
    distances (taxonomic\,+\,visual\,+\,environment) & +0.002 & [$-$0.007,\,+0.011] & 0.374 \\
    ~~visual only                                    & +0.003 & [$-$0.005,\,+0.013] & 0.257 \\
    ~~taxonomic only (full-taxonomy subset, $n{=}82$)& $-$0.006 & --- & 0.952 \\
    ~~environment only                               & $-$0.001 & [$-$0.005,\,+0.003] & 0.747 \\
    animal size (\texttt{log\_area}) [control]       & +0.005 & [$-$0.003,\,+0.014] & 0.143 \\
    \bottomrule
  \end{tabular}
\end{table}
}{}

\paragraph{Minimum detectable effect.} Table~\ref{tab:mde} quantifies the power of the primary
out-of-sample test. An effect of known standardised size $\beta$ is planted onto the real per-cell scores on
the logit scale (Equation~\ref{eq:plant}, with $z$ the standardised visual distance), whole groups are
subsampled without replacement per replicate, and the unchanged
leave-whole-group-out procedure of Section~\ref{sec:meth-cv} is asked whether it detects the plant at
$p<0.05$. The real group structure, support weights and residual noise carry through untouched, and cells at
\texttt{pDetA}=0 barely respond to the plant, so the estimates are conservative. At $\beta=0$ no replicate
is flagged, confirming the paired group-bootstrap is not anti-conservative. Generated by
src/analysis/power.py; $50$ replicates per cell.
\IfFileExists{tables/mde.tex}{% generated by src/analysis/power.py --- do not edit by hand
\begin{table}[htbp]\centering
  \caption[Minimum detectable effect]{Planted-effect power of the primary out-of-sample test
  (criterion (ii), detection, 50 replicates per cell): the fraction of
  replicates in which a species-level effect of standardised size $\beta$, planted on the real scores,
  is detected at $p<0.05$. $\Delta$MAE is the mean out-of-sample gain the plant produces. Cells at
  $\texttt{pDetA}=0$ barely respond to the plant, so these estimates are conservative.}
  \label{tab:mde}
  \footnotesize
  \begin{tabular}{lcccccccc}
    \toprule
    & \multicolumn{8}{c}{planted effect size $\beta$ (standardised logit units)} \\
    Scheme & 0.00 & 0.30 & 0.75 & 1.00 & 1.25 & 1.50 & 1.75 & 2.00 \\
    \midrule
    leave-species-out power & 0.00 & 0.00 & 0.00 & 0.00 & 0.06 & 0.58 & 0.88 & 0.98 \\
    \quad mean $\Delta$MAE at $\beta$ & +0.003 & +0.002 & +0.003 & +0.005 & +0.009 & +0.017 & +0.024 & +0.031 \\
    leave-location-out power & 0.00 & 0.00 & 0.00 & 0.08 & 0.56 & 0.98 & 1.00 & 1.00 \\
    \quad mean $\Delta$MAE at $\beta$ & +0.003 & +0.002 & +0.004 & +0.007 & +0.012 & +0.020 & +0.028 & +0.035 \\
    \bottomrule
  \end{tabular}
\end{table}
}{}

